

One way of limiting the total debt in the economy is the so-called
required reserve ratio $r$. Every bank is required by law
to set aside a fraction $r$ of the money deposited into the bank, and
this reserved money cannot be loaned further. 
If the initial amount of
money in the system (the monetary base) is $M_0$, then with loans and
borrowing the total amount of positive money available to agents is 
\begin{equation}
    M = \frac{M_0}{r},
\end{equation}
where the factor $1/r$ is called the money multiplier. 
This is how banks "create money" in the modern economy. 
They do so by creating debt obligations in the system, where the maximal debt(limited by the factor $1/r$) can be:
\begin{equation}
    M_d = M - M_0 = \frac{M_0}{r} - M_0  = M_0 \frac{1-r}{r}.
\end{equation}  




When the debt is maximized, the money distribution is again exponential but now "double sided".
Both total amounts of positive $\frac{M_0}{r}$ and negative $\frac{M_0(1-r)}{r}$ money are conserved separately
since money circulate between the agents in the
system, so there are effectively two conservation laws for
each of them.

Let $r\in(0,1)$ denote the fraction of agents in the positive sector and $1-r$ the fraction in the debt sector. The populations are initialized as
\[
N_+ = rN, \qquad N_- = (1-r)N,
\]
and the corresponding total amounts of money are
\[
M_+ = \frac{M_0}{r}, \qquad M_- = \frac{M_0(1-r)}{r}.
\]
This choice ensures that the total net wealth remains
$M_+ - M_- =  M_0$.
Hence the system has a positive money sector and a negative debt sector, while preserving the net conserved amount $M_0$.


The two pools evolve independently, and at the end the debt pool is reflected to negative values:
\[
(x_1,\dots,x_{N_+},-y_1,\dots,-y_{N_-}).
\]

This construction produces a bimodal distribution with a positive branch and a negative branch, separated by a zero-crossing at $m=0$. The parameter $r$ controls the relative weight of the two sectors. To determine the value of the reserve ratio that aligns the two branches at the origin, we impose the matching condition
\[
p_+(0)=p_-(0),
\]
where the two densities are
\[
P(m_+) = \frac{N_+}{N}\frac{1}{T_+} e^{-m_+/T_+}, \qquad m_+\ge 0,
\]
\[
P(m_-) = \frac{N_-}{N}\frac{1}{T_-} e^{m_-/T_-}, \qquad m_-\le 0.
\]

\subsection{Reserve Ratio choice of value $r$}
From section \ref{sec:debtModel}, we can see that the choice of $r$ is crucial for the model to work properly.
The choice of $r$ is based on the assumption that the total debt in the system is limited by the required reserve ratio.
To have matching populations at $m=0$, the condition $P(m_+)=P(m_-)$ reduces to
\[
\frac{N_+}{T_+} = \frac{N_-}{T_-},
\]
with
\[
T_+ = \frac{M_+}{N_+}, \qquad T_- = \frac{M_-}{N_-}.
\]
Using the definitions of the two sectors,
\[
N_+ = rN, \qquad N_- = (1-r)N,
\]
\[
M_+ = \frac{M_0}{r}, \qquad M_- = \frac{M_0(1-r)}{r},
\]
we obtain
\[
T_+ = \frac{M_0}{r^2 N}, \qquad T_- = \frac{M_0}{rN}.
\]
Substituting these expressions into the matching condition yields
\[
r \cdot \frac{r^2N}{M_0} = (1-r)\cdot \frac{rN}{M_0},
\]
which simplifies to
\[
r^2 + r - 1 = 0.
\]
Therefore the reverse ratio that makes the two distributions meet at the origin is
\[
r = \frac{\sqrt{5}-1}{2} \approx 0.618.
\]
This is the value used in the simulations to compare the positive and negative branches of the double-sided Boltzmann-Gibbs distribution.



\subsection{Simulation results}
One way of limiting the total debt in the economy, as we explained in section \ref{sec:debtModel}, is the so-called required reserve ratio $r$~\cite{xi2005required}.
The value of r was set to $r = 0.618$ in order to have the two branches of the double-sided Boltzmann-Gibbs distribution meet at the origin.
The computation of this value is shown in Appendix \ref{app:mathematics}.

The results are shown in Fig.~\ref{fig:ReverseRatioImage}, where the histogram of the money distribution is plotted along with the theoretical double-sided Boltzmann-Gibbs distribution.
\begin{figure}[H]
    \centering
    \includegraphics[width=0.4\textwidth]{ReverseRatioImage.png}
    \caption{The histogram shows the money distribution of agents after $10^7$ number of transactions, while the red and blue lines represent the theoretical double-sided Boltzmann-Gibbs distribution.
                The simulation was performed with $N=500$, $M_0 = 5000$, and $\Delta m = 1$. The reserve ratio is set to $r = 0.618$ so that the two branches of the distribution meet at the origin.}
    \label{fig:ReverseRatioImage}
\end{figure}